\documentclass[10pt,a4paper]{article}
\usepackage[margin=22mm]{geometry}
\usepackage[T1]{fontenc}
\usepackage{lmodern,microtype,amsmath,amssymb,amsthm,booktabs,array,xcolor}
\newtheorem{theorem}{Theorem}
\newtheorem{lemma}[theorem]{Lemma}
\newcommand{\bd}{\partial}
\newcommand{\card}[1]{\lvert#1\rvert}
\newcommand{\cl}{\operatorname{cl}}
\newcommand{\ZF}{Z}
\setlength{\parskip}{4pt}
\setlength{\emergencystretch}{1em}
\title{Zero forcing in the generalized Petersen graphs $P(n,4)$}
\author{Wenlin Zhang and Haobo Ma}
\date{8 October 2026 (revised note)}
\begin{document}
\maketitle
\begin{abstract}
We give a computer-assisted proof that $\ZF(P(n,4))=10$ for every
$n\ge18$, and determine the nine smaller admissible cases. The lower bound
uses the external boundary of sixteen forcing sources. Deleting a column
from a run of nine empty columns reduces the required vertex-boundary
inequality to the finite range $20\le n\le108$. The cases $n=18,19$ are
settled by forts. The finite lower bounds and the small forcing witnesses are verified by
the accompanying reproducibility package, with a manifest of SHA-256 hashes
and separately checked propositional proofs.
\end{abstract}

The graph $P(n,4)$ has vertices $u_i,v_i$, indexed by $\mathbb Z/n\mathbb Z$,
and edges $u_i u_{i+1}$, $u_i v_i$, and $v_i v_{i+4}$. We assume $n\ge9$,
so that these graphs are simple and cubic. In zero forcing, a black vertex
with exactly one white neighbor colors that neighbor black. A set is zero
forcing if this rule colors every vertex black; $\ZF(G)$ is the least size
of such a set.

\begin{theorem}\label{main}
For every $n\ge18$, $\ZF(P(n,4))=10$. For the remaining admissible values,
\[
\begin{array}{c|rrrrrrrrr}
n&9&10&11&12&13&14&15&16&17\\\hline
\ZF(P(n,4))&6&6&7&6&8&8&9&8&9.
\end{array}
\]
In particular, eighteen is the least eventual stabilization threshold.
\end{theorem}

Rashidi, Shajareh Poursalavati, and Tavakkoli \cite[Theorem 2.2]{RSPT}
proved the general upper bound $\ZF(P(n,k))\le2k+2$. Krishnan's correction
\cite{Krishnan} concerns $k=3$; its conclusion leaves the stabilization
threshold and uniform lower bound for $k\ge4$ open. It does not state a
separate conjecture for $k=4$. We use the first-forcers method from our earlier note
\emph{A uniform lower bound for zero forcing in $P(n,3)$} \cite{ZM},
by Wenlin Zhang and Haobo Ma, dated 6 October 2026 and sent to Krishnan.
That note proves his Conjecture~5, namely $\ZF(P(n,3))=8$ for every
$n\ge13$. Here we use sixteen sources in place of ten. The upper bound below is an explicit instance of the known
general construction.

\section{An explicit ten-vertex forcing set}
For $n\ge11$, start with $S=\{u_0,\ldots,u_9\}$ black. The following
sequence is legal:
\[
u_i\longrightarrow v_i\quad(1\le i\le8),\qquad
v_4\longrightarrow v_0,\qquad v_5\longrightarrow v_9.
\]
The first eight sources have both outer neighbors black. The source $v_4$
has black neighbors $u_4,v_8$ and white neighbor $v_0$; similarly $v_5$
has black neighbors $u_5,v_1$ and white neighbor $v_9$. Thus both vertices
in each of columns zero through nine are black.

For $j=10,\ldots,n-1$, perform
\[
u_{j-1}\longrightarrow u_j,\qquad
v_{j-4}\longrightarrow v_j.
\]
The induction invariant before stage $j$ is that both vertices in columns
$0,\ldots,j-1$ are black, and every vertex in columns $j,\ldots,n-1$
is still white. Thus the two sources are black and both targets are white.
In the first force, the other neighbors are $u_{j-2},v_{j-1}$;
in the second, they are $u_{j-4},v_{j-8}$. These four vertices lie in earlier
columns and are black. The first force leaves $v_j$ white, so the second
is legal too. The two forces establish the invariant for the next stage.
This colors every vertex. For $n=9,10$, taking all outer vertices gives
a zero forcing set of size $n$. Consequently $\ZF(P(n,4))\le10$ for
every admissible $n$.

\section{The lower bound from sixteen forcing sources}
For a vertex set $X$, write
\[
\bd X=\{w\notin X:\text{$w$ has a neighbor in $X$}\}
\]
for its external vertex boundary.

\begin{lemma}\label{first}
In a finite graph, if a zero forcing set $S$ performs at least $p$ forces,
the sources $X$ of its first $p$ forces satisfy
$\card X=p$ and $\card{\bd X}\le\card S$.
\end{lemma}
\begin{proof}
A vertex cannot force twice: after it forces, all its neighbors are black.
After $p$ forces the black set $B$ has size $\card S+p$, contains $X$, and
contains every neighbor of $X$. Therefore
$\bd X\subseteq B\setminus X$, whose size is $\card S$.
\end{proof}

\begin{lemma}\label{boundary}
If $n\ge20$ and $X\subseteq V(P(n,4))$ has sixteen vertices, then
$\card{\bd X}\ge10$.
\end{lemma}
\begin{proof}
Suppose $\card{\bd X}\le9$. Let $C$ be the set of occupied columns,
meaning those containing at least one vertex of $X$, and put $c=\card C$.
In each occupied column, every vertex outside $X$ belongs to $\bd X$
by its spoke. Hence
\[
2c-16\le\card{\bd X}\le9,\qquad c\le12.
\]
If $n>108$, some run of empty columns has length at least nine.
Indeed, if each of the $c$ cyclic gaps contained at most eight empty
columns, then $n\le c+8c\le108$.

Rotate the empty run to columns $0,\ldots,8$, and delete column four.
Define a bijection on the surviving column indices by
\[
f:\{0,\ldots,n-1\}\setminus\{4\}\longrightarrow\{0,\ldots,n-2\},
\qquad
f(i)=\begin{cases}i,&0\le i\le3,\\ i-1,&5\le i\le n-1.\end{cases}
\]
Extend $f$ to vertices by preserving the layer, and put $X'=f(X)$.
Every occupied index satisfies $9\le i\le n-1$. For each
$d\in\{-4,-1,1,4\}$, put $q=i+d$. Then $5\le q\le n+3$, so the
deleted column is never a neighbor index. If $q<n$, both $i$ and $q$
are above four and
\[
f(q)=q-1=f(i)+d.
\]
If $q\ge n$, its old residue $q-n$ is in $\{0,1,2,3\}$, and
\[
f(q\bmod n)=q-n=(f(i)+d)\bmod(n-1).
\]
These are all cases: $d=-1,-4$ never wraps, $d=1$ wraps only at $i=n-1$,
and $d=4$ wraps at $i=n-4,\ldots,n-1$, to old columns $0,\ldots,3$.
In particular the old edge from column $n-1$ to column three is preserved.

Conversely, an occupied image index $j=f(i)$ lies in $\{8,\ldots,n-2\}$.
For the same four offsets, $4\le j+d\le n+2$. A nonwrapping neighbor
$j+d\ge4$ lifts to old column $j+d+1=i+d$; a wrapping neighbor has
residue in $\{0,1,2,3\}$ and lifts to that same residue. Thus no new
neighbor is introduced at the deletion. This is where the full nine-column
empty buffer is used: old neighbors avoid column four, and new neighbors
have exactly the indicated lifts. Spokes are preserved as well. For each
$x\in X$, the map therefore carries its three neighbors bijectively onto
the three neighbors of $f(x)$.

Writing $N[X]=X\cup\bigcup_{x\in X}N(x)$, we obtain
\[
f(N[X])=N[X'],\qquad f(X)=X',\qquad
f(\bd X)=f(N[X]\setminus X)=N[X']\setminus X'=\bd X'.
\]
The injectivity of $f$ gives
$\card{X'}=16$ and $\card{\bd X'}=\card{\bd X}$.
Repeat until the circumference is at most 108. Starting above 108 this
procedure stops at 108, so it cannot introduce a smaller exceptional
circumference. It remains to exclude the original conditions for
$20\le n\le108$. All 89 cases are excluded by the accompanying package
(\texttt{MANIFEST.sha256}); Section~3 specifies the formulas and checks.
This completes the proof.
\end{proof}

If $n\ge20$ and $S$ had at most nine vertices, its forcing process would
have at least $2n-9\ge31$ forces. Applying Lemma~\ref{first} with $p=16$
would contradict Lemma~\ref{boundary}. Thus $\ZF(P(n,4))\ge10$ for $n\ge20$.

The boundary inequality does require a threshold: in $P(19,4)$ the set
\[
X=\{u_0,\ldots,u_9,v_3,v_4,v_5,v_6,v_{10},v_{18}\}
\]
has sixteen vertices and boundary
$\{u_{10},u_{18},v_0,v_1,v_2,v_7,v_8,v_9,v_{14}\}$, of size nine.
This explains the separate fort arguments for $n=18,19$.

\section{Finite verification}
The accompanying directory \texttt{package/} contains the complete
finite evidence for this note. Its \texttt{MANIFEST.sha256} identifies the
sources, graph data, all 100 CNFs, and compressed binary DRAT proofs by
SHA-256 hashes. The computational steps below are verified by this package.
They form part of the proof of Lemma~\ref{boundary}.

For each $20\le n\le108$, a Boolean variable $x_w$ records whether
$w\in X$. For $20\le n\le31$, a variable $z_w$ records whether $w$ is in
the closed neighborhood of $X$. For $32\le n\le108$, $z_w$ instead records
membership in $\bd X$. The respective equivalences are
\[
z_w\ \longleftrightarrow\ x_w\lor\bigvee_{t\sim w}x_t,
\qquad\text{or}\qquad
z_w\ \longleftrightarrow\ \neg x_w\land\bigvee_{t\sim w}x_t.
\]
The cardinality conditions are $\sum_w x_w=16$ and, respectively,
$\sum_w z_w\le25$ or $\sum_w z_w\le9$. They are equivalent to the
boundary conditions because the closed neighborhood is the disjoint
union of $X$ and $\bd X$. We also impose $x_{u_0}\lor x_{v_0}$, which is
valid after rotating any nonempty $X$. The equivalences are expanded
directly into clauses, and the cardinalities use sequential counters
\cite{Sinz}. The accompanying package verifies that every formula is unsatisfiable:
it regenerates and byte-compares each CNF, audits its semantics, and checks
the corresponding binary DRAT proof with DRAT-trim \cite{DRAT}.

For the smaller zero forcing numbers, a nonempty set $F$ is a \emph{fort}
if every vertex outside $F$ has zero or at least two neighbors in $F$.
Every zero forcing set meets $F$: the first force into a disjoint fort
would have exactly one neighbor in it, contradicting this condition.
For each $9\le n\le19$, finite fort families exclude seed sets of size
at most one less than the asserted value.

Here the Boolean variables $s_w$ record the seeds, and every fort adds
the clause $\bigvee_{w\in F}s_w$. The simultaneous column rotations $i\mapsto i+a$ and reflections
$i\mapsto a-i$ on both layers are automorphisms: they preserve spokes and
send the same-layer differences $\pm1$ and $\pm4$ to themselves.
Under this group the directed edges have four orbits: outer edges,
inner edges, spokes directed from outer to inner, and spokes directed
from inner to outer. Rotating the source to column zero and reflecting
when necessary gives the four representatives
\[
u_0\to v_0,\quad u_0\to u_1,\quad
v_0\to u_0,\quad v_0\to v_4.
\]
For each case, the source and its other two neighbors must be seeds,
and the target must not be a seed. The formula requires at least one
case and at most the specified number of seeds. A zero forcing set of size at most the specified bound is not the full
vertex set; hence it must admit an initial force, since otherwise the
process cannot start. Normalizing that directed edge gives one of the
four cases. Thus every candidate zero forcing set satisfies a normalized
formula. The package checks every listed set directly against the fort
condition and checks all eleven UNSAT proofs. This proves the finite lower
bounds, including ten at $n=18,19$.

For $9\le n\le17$, the following seed sets attain the stated values.
Complete, ordered forcing sequences are printed in Appendix~\ref{sequences}
and verified by the package.
\[
\begin{array}{c|l}
n&\text{zero forcing set}\\\hline
9&u_0,u_1,u_2,u_3,u_8,v_4\\
10&u_0,u_1,u_2,u_3,u_9,v_4\\
11&u_0,\ldots,u_5,u_{10}\\
12&u_0,u_1,u_2,u_{11},v_4,v_5\\
13&u_0,u_1,u_2,u_3,u_5,u_6,u_7,u_{12}\\
14&u_0,\ldots,u_5,u_8,u_{13}\\
15&u_0,\ldots,u_7,u_{14}\\
16&u_0,\ldots,u_6,u_{15}\\
17&u_0,\ldots,u_7,u_{16}.
\end{array}
\]
Together with these witnesses and the uniform ten-vertex construction,
the verified lower bounds prove Theorem~\ref{main}.

\paragraph{Reproducing the finite proof.}
The file \texttt{generate.py} deterministically emits the boundary formulas
and the fort formulas from the stored vertex lists and clause-order data.
The separate \texttt{semantic\_check.py} rebuilds $P(n,4)$ from the direct
neighbor formulas. It validates every neighborhood clause, every sequential
counter grid, the rotation clause, the four first-force cases, and every
fort-hitting clause; it also checks that each fort is nonempty and satisfies
the fort condition. It checks each forcing step against the current black
set, and checks the buffered-deletion identities. This checker imports
neither the generator nor the SAT encoding library.

The original proofs were generated by CaDiCaL~1.9.5 through
\texttt{python-sat}~1.9.dev15. The package includes that pinned distribution's
source archive and the producing-run generator sources. DRAT-trim is pinned
to source revision \texttt{2e3b2dc} (the complete revision and source hash
are recorded in \texttt{VERSIONS.json}); its source and license are included.
Verification needs Python~3.11 or later with its standard library and a C
compiler, and runs offline. From the accompanying \texttt{package/} directory,
execute
\begin{quote}
\verb|./verify_all --log-dir ../verification-logs|
\end{quote}
This command checks the manifest, regenerates all CNFs, compares their bytes
with the archived files, performs the semantic and forcing checks, builds
DRAT-trim using \texttt{cc -O2}, decompresses each proof, checks its raw hash,
and runs \texttt{drat-trim CNF PROOF -i -t 600}. Each of the 100 checks must
print \texttt{s VERIFIED}; the driver then prints \texttt{PASS ALL}.
The recorded complete run is supplied alongside the package. Exact setup
and fresh proof-generation commands are in \texttt{README.md}. The smaller
\texttt{package-lite/} omits DRAT files and includes their regeneration
commands; its semantic-only check does not establish UNSAT.

\appendix
\section{Explicit small-case forcing sequences}\label{sequences}
For each row below, start from the seed set in Section~3 and perform the
arrows in the printed order, reading left to right and then top to bottom.
Every source is black and has exactly the indicated single white neighbor.
The package stores the same sequences with vertex codes $u_i=i$, $v_i=n+i$.


\noindent\textbf{$n=9$.}\par\nopagebreak\[
\begin{gathered}
u_{0}\to v_{0}\quad u_{1}\to v_{1}\quad u_{2}\to v_{2}\quad v_{0}\to v_{5}\quad v_{1}\to v_{6}\\
v_{2}\to v_{7}\quad v_{5}\to u_{5}\quad v_{6}\to u_{6}\quad u_{5}\to u_{4}\quad u_{6}\to u_{7}\\
u_{8}\to v_{8}\quad v_{7}\to v_{3}
\end{gathered}\]
\noindent\textbf{$n=10$.}\par\nopagebreak\[
\begin{gathered}
u_{0}\to v_{0}\quad u_{1}\to v_{1}\quad u_{2}\to v_{2}\quad v_{0}\to v_{6}\quad v_{2}\to v_{8}\\
v_{4}\to u_{4}\quad v_{6}\to u_{6}\quad v_{8}\to u_{8}\quad u_{3}\to v_{3}\quad u_{4}\to u_{5}\\
u_{5}\to v_{5}\quad u_{6}\to u_{7}\quad u_{7}\to v_{7}\quad u_{9}\to v_{9}
\end{gathered}\]
\noindent\textbf{$n=11$.}\par\nopagebreak\[
\begin{gathered}
u_{0}\to v_{0}\quad u_{1}\to v_{1}\quad u_{2}\to v_{2}\quad u_{3}\to v_{3}\quad u_{4}\to v_{4}\\
v_{0}\to v_{7}\quad v_{3}\to v_{10}\quad v_{4}\to v_{8}\quad v_{7}\to u_{7}\quad v_{8}\to u_{8}\\
v_{10}\to v_{6}\quad u_{7}\to u_{6}\quad u_{8}\to u_{9}\quad u_{9}\to v_{9}\quad v_{1}\to v_{5}
\end{gathered}\]
\noindent\textbf{$n=12$.}\par\nopagebreak\[
\begin{gathered}
u_{0}\to v_{0}\quad u_{1}\to v_{1}\quad v_{0}\to v_{8}\quad v_{1}\to v_{9}\quad v_{4}\to u_{4}\\
v_{5}\to u_{5}\quad v_{8}\to u_{8}\quad v_{9}\to u_{9}\quad u_{4}\to u_{3}\quad u_{5}\to u_{6}\\
u_{8}\to u_{7}\quad u_{9}\to u_{10}\quad u_{10}\to v_{10}\quad u_{11}\to v_{11}\quad u_{2}\to v_{2}\\
u_{3}\to v_{3}\quad u_{6}\to v_{6}\quad u_{7}\to v_{7}
\end{gathered}\]
\noindent\textbf{$n=13$.}\par\nopagebreak\[
\begin{gathered}
u_{0}\to v_{0}\quad u_{1}\to v_{1}\quad u_{2}\to v_{2}\quad u_{6}\to v_{6}\quad v_{2}\to v_{11}\\
v_{6}\to v_{10}\quad v_{10}\to u_{10}\quad v_{1}\to v_{5}\quad v_{5}\to v_{9}\quad v_{9}\to u_{9}\\
u_{5}\to u_{4}\quad u_{9}\to u_{8}\quad u_{10}\to u_{11}\quad u_{12}\to v_{12}\quad v_{0}\to v_{4}\\
v_{4}\to v_{8}\quad v_{11}\to v_{7}\quad v_{12}\to v_{3}
\end{gathered}\]
\noindent\textbf{$n=14$.}\par\nopagebreak\[
\begin{gathered}
u_{0}\to v_{0}\quad u_{1}\to v_{1}\quad u_{2}\to v_{2}\quad u_{3}\to v_{3}\quad u_{4}\to v_{4}\\
v_{0}\to v_{10}\quad v_{4}\to v_{8}\quad v_{8}\to v_{12}\quad v_{12}\to u_{12}\quad u_{12}\to u_{11}\\
u_{13}\to v_{13}\quad v_{2}\to v_{6}\quad v_{3}\to v_{7}\quad v_{6}\to u_{6}\quad v_{10}\to u_{10}\\
v_{13}\to v_{9}\quad u_{5}\to v_{5}\quad u_{6}\to u_{7}\quad u_{8}\to u_{9}\quad u_{11}\to v_{11}
\end{gathered}\]
\noindent\textbf{$n=15$.}\par\nopagebreak\[
\begin{gathered}
u_{0}\to v_{0}\quad u_{1}\to v_{1}\quad u_{2}\to v_{2}\quad u_{3}\to v_{3}\quad u_{4}\to v_{4}\\
u_{5}\to v_{5}\quad u_{6}\to v_{6}\quad v_{0}\to v_{11}\quad v_{1}\to v_{12}\quad v_{2}\to v_{13}\\
v_{4}\to v_{8}\quad v_{5}\to v_{9}\quad v_{6}\to v_{10}\quad v_{8}\to u_{8}\quad v_{9}\to u_{9}\\
v_{12}\to u_{12}\quad v_{13}\to u_{13}\quad u_{7}\to v_{7}\quad u_{9}\to u_{10}\quad u_{10}\to u_{11}\\
u_{14}\to v_{14}
\end{gathered}\]
\noindent\textbf{$n=16$.}\par\nopagebreak\[
\begin{gathered}
u_{0}\to v_{0}\quad u_{1}\to v_{1}\quad u_{2}\to v_{2}\quad u_{3}\to v_{3}\quad u_{4}\to v_{4}\\
u_{5}\to v_{5}\quad v_{0}\to v_{12}\quad v_{1}\to v_{13}\quad v_{4}\to v_{8}\quad v_{5}\to v_{9}\\
v_{8}\to u_{8}\quad v_{9}\to u_{9}\quad v_{12}\to u_{12}\quad v_{13}\to u_{13}\quad u_{8}\to u_{7}\\
u_{9}\to u_{10}\quad u_{12}\to u_{11}\quad u_{13}\to u_{14}\quad u_{14}\to v_{14}\quad u_{15}\to v_{15}\\
v_{2}\to v_{6}\quad v_{3}\to v_{7}\quad v_{6}\to v_{10}\quad v_{7}\to v_{11}
\end{gathered}\]
\noindent\textbf{$n=17$.}\par\nopagebreak\[
\begin{gathered}
u_{0}\to v_{0}\quad u_{1}\to v_{1}\quad u_{2}\to v_{2}\quad u_{3}\to v_{3}\quad u_{4}\to v_{4}\\
u_{5}\to v_{5}\quad u_{6}\to v_{6}\quad v_{0}\to v_{13}\quad v_{1}\to v_{14}\quad v_{2}\to v_{15}\\
v_{4}\to v_{8}\quad v_{5}\to v_{9}\quad v_{6}\to v_{10}\quad v_{9}\to u_{9}\quad v_{10}\to u_{10}\\
v_{13}\to u_{13}\quad v_{14}\to u_{14}\quad u_{9}\to u_{8}\quad u_{10}\to u_{11}\quad u_{13}\to u_{12}\\
u_{14}\to u_{15}\quad u_{16}\to v_{16}\quad v_{3}\to v_{7}\quad v_{7}\to v_{11}\quad v_{8}\to v_{12}
\end{gathered}\]

\begin{thebibliography}{9}
\bibitem{RSPT}
S.~Rashidi, N.~Shajareh Poursalavati, and M.~Tavakkoli,
\emph{Computing the zero forcing number for generalized Petersen graphs},
J.~Algebra Comb.~Discrete Struct.~Appl.~7 (2020), no.~2, 183--193.
doi:10.13069/jacodesmath.729465.

\bibitem{Krishnan}
A.~Krishnan, \emph{A correction to the Zero Forcing Number of the
Generalized Petersen Graphs $P(n,3)$},
arXiv:2607.19412v1 (2026),
Introduction and Section~5.

\bibitem{ZM}
W.~Zhang and H.~Ma, \emph{A uniform lower bound for zero forcing in
$P(n,3)$}, mathematical note, 6 October 2026; communicated to A.~Krishnan.
Proves Conjecture~5 of \cite{Krishnan}: $\ZF(P(n,3))=8$ for $n\ge13$.

\bibitem{Sinz}
C.~Sinz, \emph{Towards an optimal CNF encoding of Boolean cardinality
constraints}, Principles and Practice of Constraint Programming,
LNCS 3709 (2005), 827--831.
doi:10.1007/11564751\_73.

\bibitem{DRAT}
M.~Heule and contributors, \emph{DRAT-trim},
Source revision and SHA-256 hash in the accompanying \texttt{VERSIONS.json}.
The source and its license are included with the supporting files.
\end{thebibliography}
\end{document}
